\documentclass[conference,a4paper]{IEEEtran}
\IEEEoverridecommandlockouts

\usepackage{cite}
\usepackage{amsmath,amssymb,amsfonts}
\usepackage{graphicx}
\usepackage{booktabs}
\usepackage{multirow}
\usepackage{xcolor}
\usepackage{url}
\usepackage{placeins}
\usepackage[hidelinks]{hyperref}

% ---------- Float placement (keeps figures before the references) ----------
\renewcommand{\topfraction}{0.95}
\renewcommand{\bottomfraction}{0.9}
\renewcommand{\textfraction}{0.05}
\renewcommand{\floatpagefraction}{0.8}
\renewcommand{\dbltopfraction}{0.95}
\renewcommand{\dblfloatpagefraction}{0.8}
\setcounter{topnumber}{4}
\setcounter{dbltopnumber}{4}

% ---------- Guidance text (switch off before submission) ----------
\newif\ifshowguidance
\showguidancefalse          % change to \showguidancefalse to hide guidance
\newcommand{\guide}[1]{\ifshowguidance{\color{gray}\itshape [#1]}\fi}

% ---------- Figure helper: real image if present, placeholder if not ----------
% \projfig{file name in figures/}{width}{placeholder text}
\newlength{\projw}\newlength{\projh}
\newcommand{\projfig}[3]{%
  \IfFileExists{figures/#1}%
    {\includegraphics[width=#2]{figures/#1}}%
    {\setlength{\projw}{#2}\setlength{\projh}{0.6\projw}%
     \fbox{\parbox[c][\projh][c]{\dimexpr\projw-2\fboxsep-2\fboxrule\relax}{%
       \centering\footnotesize\color{gray}Placeholder: #3\\[2pt]%
       \texttt{\detokenize{#1}}}}}}

\begin{document}

% =====================================================================
\title{Handwritten Digit Recognition Using Linear and Kernel Support Vector Machines\\
{\large Course Project -- Machine Learning, M.Tech (Artificial Intelligence), Semester I}}

\author{\IEEEauthorblockN{Anamika J}
\IEEEauthorblockA{\textit{Register No.: TVE26CSAI03}\\
\textit{Department of Computer Science, College of Engineering, Trivandrum}\\
{Thiruvananthapuram}, India\\
{csaip260050@cet.ac.in}}}

\maketitle

% =====================================================================
\begin{abstract}
This project studies support vector machines (SVMs) for binary classification, from two-dimensional toy problems
to the recognition of handwritten digits. A soft-margin linear SVM was trained on two linearly separable
two-dimensional data sets, and the effect of the regularisation term $C$ on the margin and on the number of
support vectors was measured. The same model was then used to separate the digits 0 and 1 (12665 training and
2115 test images). A hard-margin SVM classified every training image correctly and made two test errors (0.095\%).
Evaluating ten values of $C$ showed that a slightly softer margin can lower the test error to one image (0.047\%),
and five-fold cross-validation on the training set found such a value without touching the test data. Finally,
a radial basis function (RBF) kernel SVM was applied to a data set that no straight line can separate (a linear SVM
misclassifies 7.1\% of its points). Increasing the kernel parameter from 1 to 1000 turned the decision boundary
from almost straight and underfitting into a fragmented boundary that memorises the training data. The results
illustrate how $C$ and the kernel width control the complexity of an SVM.
\end{abstract}

\begin{IEEEkeywords}
Support vector machine, soft margin, regularisation, RBF kernel, handwritten digit recognition, cross-validation
\end{IEEEkeywords}

% =====================================================================
\section{Introduction}
Automatic recognition of handwritten digits is used in practice to read postal codes, bank cheques and
forms. The task is hard because each digit can be written in many shapes, thicknesses and slants, so a
classifier must separate the classes and still tolerate unusual writing.

Support vector machines~\cite{cortes1995,boser1992} are maximum-margin classifiers: among all hyperplanes that
separate two classes, an SVM chooses the one with the largest distance to the closest training points. A soft-margin
formulation tolerates some margin violations, and kernel functions extend the method to non-linear boundaries~\cite{scholkopf2002}.

This project has three objectives: (1) to train a soft-margin linear SVM on two two-dimensional data sets and to
study the effect of the regularisation term $C$; (2) to apply a linear SVM, with and without regularisation, to separate
handwritten 0s and 1s and to select $C$ properly; and (3) to train an RBF kernel SVM on a data set that is not linearly
separable and to study the effect of the kernel parameter $\gamma$.
Section~II reviews the SVM formulation, Section~III describes the data, and Section~IV the implementation.
The results are given in Section~V, discussed in Section~VI, and Section~VII concludes the paper.

% =====================================================================
\section{Background}
Let $\{(\mathbf{x}_i,y_i)\}_{i=1}^{n}$ be the training data with $y_i\in\{-1,+1\}$. The soft-margin SVM solves
\begin{equation}
\begin{aligned}
\min_{\mathbf{w},b,\boldsymbol{\xi}}\ & \tfrac12\|\mathbf{w}\|^2 + C\sum_{i=1}^{n}\xi_i\\
\text{s.t.}\ & y_i(\mathbf{w}^{T}\mathbf{x}_i+b)\ge 1-\xi_i,\ \ \xi_i\ge 0 .
\end{aligned}
\label{eq:primal}
\end{equation}
The margin has width $2/\|\mathbf{w}\|$, the slack variable $\xi_i$ measures how far point $i$ violates it, and
$C>0$ balances a wide margin against few violations. Setting all $\xi_i=0$ (equivalently $C\to\infty$) gives the
hard-margin SVM, which exists only for linearly separable data.

The dual problem is
\begin{equation}
\begin{aligned}
\max_{\boldsymbol{\alpha}}\ & \sum_i\alpha_i-\tfrac12\sum_{i,j}\alpha_i\alpha_j y_i y_j K(\mathbf{x}_i,\mathbf{x}_j)\\
\text{s.t.}\ & 0\le\alpha_i\le C,\ \ \sum_i\alpha_i y_i=0,
\end{aligned}
\label{eq:dual}
\end{equation}
and the classifier is $f(\mathbf{x})=\sum_i\alpha_i y_i K(\mathbf{x}_i,\mathbf{x})+b$. Only the support vectors
($\alpha_i>0$) contribute. For the linear kernel $K(\mathbf{x}_i,\mathbf{x}_j)=\mathbf{x}_i^{T}\mathbf{x}_j$ and
$\mathbf{w}=\sum_i\alpha_i y_i\mathbf{x}_i$. If $C$ is larger than every $\alpha_i$ of the hard-margin solution,
the constraint $\alpha_i\le C$ is never active and the soft-margin solution coincides with the hard-margin one.

The RBF kernel
\begin{equation}
K(\mathbf{x}_i,\mathbf{x}_j)=\exp\!\big(-\gamma\|\mathbf{x}_i-\mathbf{x}_j\|^2\big),\qquad\gamma>0,
\label{eq:rbf}
\end{equation}
produces non-linear boundaries. The quantity $1/\sqrt{\gamma}$ acts as the length scale over which a support vector
influences the decision.

% =====================================================================
\section{Datasets and Preprocessing}
\begin{table}[!t]
\caption{Summary of Data Sets}
\label{tab:data}
\centering
\begin{tabular}{@{}lccc@{}}
\toprule
Data set & Features & Train size & Test size\\
\midrule
Data set 1 & 2 & 1500 & 500\\
Data set 2 & 2 & 1500 & 500\\
Digits 0/1 (images) & 784 & 12665 & 2115\\
Data set 3 & 2 & 211 & --\\
\bottomrule
\end{tabular}
\end{table}

Table~\ref{tab:data} lists the four data sets.

\emph{Data sets 1 and 2.} Each row holds two real-valued features (roughly 35 to 167) and a label $\pm1$. The
training sets contain 748/752 and 729/771 positive/negative points, and the test sets 252/248 and 271/229.
In data set 1 the two classes are far apart, whereas in data set 2 they nearly touch.

\emph{Digits.} Each line of the files stores one $28\times28$ grey-scale image in sparse form,
\texttt{label index:value \dots}, with 1-based pixel indices and only the non-zero pixels listed. Label $+1$ is
the digit 1 and $-1$ is the digit 0. The training set has 6742 ones and 5923 zeros; the test set has 1135 ones and
980 zeros. These are the 0 and 1 classes of the MNIST database~\cite{lecun1998}. Every image was converted to a
dense 784-dimensional vector and the grey values were scaled from $[0,255]$ to $[0,1]$.

\emph{Data set 3.} It has 211 two-dimensional points (106 positive, 105 negative) with coordinates between about $-0.66$ and
$0.57$, and is used only for training.

% =====================================================================
\section{Methodology}
All experiments were run in Google Colab with Python, NumPy and Matplotlib. The SVMs were trained with the
\texttt{SVC} class of scikit-learn~\cite{pedregosa2011}, which wraps the LIBSVM solver~\cite{chang2011} for the dual
problem~\eqref{eq:dual}.

\emph{Soft-margin linear SVM (Part~1).} \texttt{SVC(kernel="linear", C=$C$)} was fitted to the training set of each
data set. The decision boundary $f(\mathbf{x})=0$ and the margins $f(\mathbf{x})=\pm1$ were drawn by evaluating
$f$ on a dense grid. The test error is the fraction of misclassified test points. To study $C$, ten values from
$10^{-6}$ to $10^{3}$ were used, and for each the training and test errors, the number of support vectors and the margin
$2/\|\mathbf{w}\|$ were recorded.

\emph{Hard-margin SVM (Part~2(i)).} A hard margin (all $\xi_i=0$) was approximated by $C=10^{6}$, far above the
largest dual variable, which was checked afterwards. Because this model makes no training errors, a training failure
was illustrated with a strongly regularised model ($C=10^{-4}$), and one image it misclassifies was analysed.
The weight vector $\mathbf{w}$ was displayed as a $28\times28$ image, and the training images with the smallest
functional margin $y_i f(\mathbf{x}_i)$ were inspected.

\emph{Choosing $C$ (Part~2(ii)).} Ten values of $C$ on an approximately logarithmic scale between $10^{-4}$ and $10$ were
tried. The upper end lies far above the largest dual variable of the hard-margin solution, where nothing changes, and
below $10^{-4}$ errors increase visibly. To choose $C$ without using the test set, stratified five-fold cross-validation
was run on the training set, and the model with the lowest cross-validation error was then evaluated once on the test set.

\emph{RBF kernel SVM (Part~3).} \texttt{SVC(kernel="rbf", gamma=$\gamma$, C=1)} was trained on data set 3 for
$\gamma\in\{1,10,100,1000\}$. A linear SVM with $C=1000$ was also trained to test whether a linear boundary exists.
Error is always the fraction of misclassified examples.

% =====================================================================
\section{Experiments and Results}

% ---------------------------------------------------------------------
\subsection{Part 1: Soft-Margin SVM on Two-Dimensional Data}
Fig.~\ref{fig:ds_train} shows the boundary for $C=1$. For data set 1 the SVM found
$\mathbf{w}=(-0.125,\,0.127)$, $b=-0.298$, a margin of 11.24 and 3 support vectors. For data set 2 it found
$\mathbf{w}=(-1.056,\,1.077)$, $b=-2.762$, a margin of only 1.33 and 5 support vectors. Both models classify all
training points correctly and misclassify none of the 500 test points of their data set (test error 0, Fig.~\ref{fig:ds_test}).

\begin{figure*}[!t]
\centering
\projfig{fig_ds_train.png}{0.78\textwidth}{training data with decision boundary, margins and support vectors}
\caption{Linear soft-margin SVM ($C=1$) on the training data of data set 1 (left) and data set 2 (right). The solid
line is $f(\mathbf{x})=0$, the dashed lines are $f(\mathbf{x})=\pm1$, and circled points are support vectors.}
\label{fig:ds_train}
\end{figure*}

\begin{figure*}[!t]
\centering
\projfig{fig_ds_test.png}{0.78\textwidth}{test data with decision boundary}
\caption{Test points of data set 1 (left) and data set 2 (right) with the boundary learned from the training data
($C=1$). No test point is misclassified.}
\label{fig:ds_test}
\end{figure*}

Ten values of $C$ were tried, from $10^{-6}$ to $10^{3}$ (a wide range chosen so that both under- and over-regularised models appear).
Table~\ref{tab:C12} lists the errors, support vectors and margins; Figs.~\ref{fig:Call1} and~\ref{fig:Call2} show the boundary for every
value, and Fig.~\ref{fig:errC12} plots the error, the number of support vectors and the margin against $C$. Both sets are linearly separable, and for
$C\ge10^{-4}$ the training and test errors are zero on both. A small $C$ makes violations cheap: the margin becomes wide and
most points turn into support vectors (1274 and 1450 at $C=10^{-6}$). A large $C$ approaches the hard-margin SVM with a narrow margin
and only 3 support vectors. For data set 1 the solution stops changing at $C\approx0.1$ (margin 11.24); for data set 2 the margin
keeps shrinking until $C\approx10$ (margin 1.18), because its classes are much closer. If $C$ is too small the model underfits: at
$C=10^{-6}$ the errors on data set 2 rise to 8.4\% (training) and 9.8\% (test).

\begin{table*}[!t]
\caption{Effect of $C$ on the Linear SVM for Data Sets 1 and 2 (SV = number of support vectors, margin $=2/\|\mathbf{w}\|$)}
\label{tab:C12}
\centering
\begin{tabular}{@{}r cccc cccc@{}}
\toprule
 & \multicolumn{4}{c}{Data set 1} & \multicolumn{4}{c}{Data set 2}\\
\cmidrule(lr){2-5}\cmidrule(lr){6-9}
$C$ & Train err. & Test err. & SV & Margin & Train err. & Test err. & SV & Margin\\
\midrule
$10^{-6}$ & 0.0000 & 0.0000 & 1274 & 73.28 & 0.0840 & 0.0980 & 1450 & 85.38\\
$10^{-5}$ & 0.0000 & 0.0000 & 444 & 35.47 & 0.0007 & 0.0000 & 750 & 30.30\\
$10^{-4}$ & 0.0000 & 0.0000 & 116 & 20.71 & 0.0000 & 0.0000 & 284 & 15.22\\
$10^{-3}$ & 0.0000 & 0.0000 & 24 & 14.26 & 0.0000 & 0.0000 & 104 & 7.59\\
$10^{-2}$ & 0.0000 & 0.0000 & 4 & 11.56 & 0.0000 & 0.0000 & 35 & 4.21\\
$10^{-1}$ & 0.0000 & 0.0000 & 3 & 11.24 & 0.0000 & 0.0000 & 10 & 2.41\\
$1$ & 0.0000 & 0.0000 & 3 & 11.24 & 0.0000 & 0.0000 & 5 & 1.33\\
$10$ & 0.0000 & 0.0000 & 3 & 11.24 & 0.0000 & 0.0000 & 3 & 1.18\\
$10^{2}$ & 0.0000 & 0.0000 & 3 & 11.24 & 0.0000 & 0.0000 & 3 & 1.18\\
$10^{3}$ & 0.0000 & 0.0000 & 3 & 11.24 & 0.0000 & 0.0000 & 3 & 1.18\\
\bottomrule
\end{tabular}
\end{table*}

\begin{figure*}[!t]
\centering
\projfig{fig_C_all_ds1.png}{0.95\textwidth}{data set 1, ten values of C}
\caption{Data set 1: decision boundary (solid) and margins (dashed) of the linear SVM for ten values of $C$ from $10^{-6}$ to $10^{3}$.
The number of support vectors (SV) is given in each title.}
\label{fig:Call1}
\end{figure*}

\begin{figure*}[!t]
\centering
\projfig{fig_C_all_ds2.png}{0.95\textwidth}{data set 2, ten values of C}
\caption{Data set 2: the same study as Fig.~\ref{fig:Call1}. The classes are much closer, so the margin keeps shrinking up to $C\approx10$.}
\label{fig:Call2}
\end{figure*}

\begin{figure*}[!t]
\centering
\projfig{fig_err_C_ds.png}{0.88\textwidth}{error, support vectors and margin versus C}
\caption{Effect of $C$ on data set 1 (top row) and data set 2 (bottom row): training and test error (left), number of support vectors
(middle) and margin width $2/\|\mathbf{w}\|$ (right).}
\label{fig:errC12}
\end{figure*}

% ---------------------------------------------------------------------
\subsection{Part 2(i): Hard-Margin SVM on Digits 0 and 1}
\begin{figure}[!t]
\centering
\projfig{fig_digit_samples.png}{\columnwidth}{sample digits and class averages}
\caption{Training samples: six 0s (top row), six 1s (bottom row), and the average 0 and average 1 of the training set.}
\label{fig:samples}
\end{figure}

Samples of the data are shown in Fig.~\ref{fig:samples}. The hard-margin SVM converged with 84 support vectors; the largest dual
variable is $\alpha_{\max}=0.395$, and the smallest functional margin over the training set is $\min_i y_if(\mathbf{x}_i)=0.9995\approx1$. The results are:
\begin{itemize}
\item training error $=0/12665=0$;
\item test error $=2/2115=0.00095$ (0.095\%).
\end{itemize}
The training error is zero because the 0/1 training set is linearly separable in the 784-dimensional pixel space, so
there is no misclassified training image to pick for the hard-margin model. To study a training failure, a strongly regularised
model with $C=10^{-4}$ was used instead; it misclassifies 30 training images, and the first of them, training image \#56
(Fig.~\ref{fig:mis}), was analysed.

Image \#56 is a 1 written in the style with a hooked head and a long foot stroke. It has 165 inked pixels, while an
average 1 has about 86 and an average 0 about 192, and its ink is spread over a wide area. A linear SVM only forms a weighted sum of
pixel intensities, so such an image gets a decision value of $f(\mathbf{x})=-0.05$ for $C=10^{-4}$ and is assigned to class 0. The hard-margin SVM
classifies it correctly, but only as a support vector with $yf(\mathbf{x})=1.0$. The two test errors of the hard-margin SVM show the same
weakness: \#1664 is a very thick, slanted 1 and \#2031 is an open, broken 0 with a stray blob.

\begin{figure*}[!t]
\centering
\projfig{fig_misclassified.png}{0.9\textwidth}{misclassified images}
\caption{From left: training image \#56, a 1 misclassified as a 0 by the model with $C=10^{-4}$; the average 0 and average 1
of the training set; and the two test images misclassified by the hard-margin SVM.}
\label{fig:mis}
\end{figure*}

Fig.~\ref{fig:weights} explains why. The weight image is mainly positive (red) near the centre, where a 1 has its ink, and mostly
negative (blue) around it, where a 0 has its strokes. The training images lying exactly on the margin are the atypical ones: a 0 with an
inward hook, a very slanted thin 1 with a stray dot, a 1 with a doubled stroke and a 1 with a heavy foot and head.

\begin{figure*}[!t]
\centering
\projfig{fig_weights.png}{0.9\textwidth}{weights and margin images}
\caption{Left: weight vector $\mathbf{w}$ of the hard-margin SVM as a $28\times28$ image (red pushes towards digit 1, blue
towards digit 0). Right: the four training images with the smallest $y_if(\mathbf{x}_i)$; all are support vectors on the margin.}
\label{fig:weights}
\end{figure*}

% ---------------------------------------------------------------------
\subsection{Part 2(ii): Choosing the Regularisation Term $C$}
\emph{Range of $C$.} The largest dual variable of the hard-margin solution is about 0.4, so for every $C$ above this value the box constraint
$\alpha_i\le C$ is inactive and the hard-margin solution is recovered. Errors start to appear at about $C=10^{-4}$. Ten values on an
approximately logarithmic scale between $10^{-4}$ and $10$ were therefore evaluated.

The results are in Table~\ref{tab:digitsC}, which also gives the five-fold cross-validation error, and in Figs.~\ref{fig:digitsC}
and~\ref{fig:weightsC}. The training error
falls steadily as $C$ grows, from 30 mistakes at $C=10^{-4}$ to 2 at $C=0.1$, and is zero from $C=0.3$ on. The test error is not
monotone.

\emph{A. Best training error.} The training error first reaches its minimum of zero at $C=0.3$ and stays zero for every larger $C$.

\emph{B. Comparison with Part 2(i).} At $C=0.3$ the test error is $2/2115=0.00095$, exactly equal to that of the hard-margin SVM. Since
$0.3$ is close to $\alpha_{\max}$ and every larger value gives the hard-margin solution, the test predictions do not change.

\emph{C. Optimising $C$ for the test error.} The lowest test error, $1/2115=0.00047$, occurs at $C=0.01$ and $C=0.03$, although the training
error there is not zero: a slightly softer margin ignores a few unusual training images and generalises better. Choosing $C$ by looking at the test
error, however, tunes on the test set and gives an optimistic estimate. The proper procedure is cross-validation on the training data.
Five-fold cross-validation selects $C=0.03$ (CV error 0.00095), and this model makes one test error (0.047\%).

\begin{table}[!t]
\caption{Linear SVM on Digits 0/1 for Different $C$ (number of misclassified images in parentheses)}
\label{tab:digitsC}
\centering
\small
\setlength{\tabcolsep}{3pt}
\begin{tabular}{@{}rcccc@{}}
\toprule
$C$ & Train error & Test error & SV & 5-fold CV\\
\midrule
$10^{-4}$ & 0.00237 (30) & 0.00189 (4) & 1808 & 0.00237\\
$3\times10^{-4}$ & 0.00197 (25) & 0.00095 (2) & 897 & 0.00205\\
$10^{-3}$ & 0.00142 (18) & 0.00095 (2) & 436 & 0.00174\\
$3\times10^{-3}$ & 0.00118 (15) & 0.00095 (2) & 240 & 0.00150\\
$10^{-2}$ & 0.00071 (9) & 0.00047 (1) & 151 & 0.00111\\
$3\times10^{-2}$ & 0.00047 (6) & 0.00047 (1) & 103 & \textbf{0.00095}\\
$10^{-1}$ & 0.00016 (2) & 0.00095 (2) & 89 & 0.00134\\
$0.3$ & 0.00000 (0) & 0.00095 (2) & 85 & 0.00111\\
$1$ & 0.00000 (0) & 0.00095 (2) & 84 & 0.00111\\
$10$ & 0.00000 (0) & 0.00095 (2) & 84 & 0.00111\\
\bottomrule
\end{tabular}
\end{table}

\begin{figure*}[!t]
\centering
\projfig{fig_digits_C.png}{0.92\textwidth}{error versus C for the digits}
\caption{Linear SVM on digits 0/1 for ten values of $C$. Left: training, test and five-fold cross-validation error (\%); the dotted line is the
test error of the hard-margin SVM. Middle: number of misclassified training and test images. Right: number of support vectors.}
\label{fig:digitsC}
\end{figure*}

\begin{figure}[!t]
\centering
\projfig{fig_weights_C.png}{\columnwidth}{weight images for different C}
\caption{Weight vector $\mathbf{w}$ for $C=10^{-4},10^{-2},10^{-1}$ and $1$ (red pushes towards digit 1, blue towards digit 0). For small $C$ the
weights are smooth and dominated by the average class shapes; for larger $C$ they become noisy and are shaped by individual hard examples.}
\label{fig:weightsC}
\end{figure}

% ---------------------------------------------------------------------
\subsection{Part 3: RBF Kernel SVM}
Fig.~\ref{fig:ds3} shows data set 3. A straight line cannot separate the positives from the negatives: even a linear SVM with
$C=1000$, which tries hard to fit every point, misclassifies 15 of the 211 points (error 0.0711). Hence \textbf{no linear decision
boundary exists} for this data set.

\begin{figure}[!t]
\centering
\projfig{fig_ds3.png}{0.8\columnwidth}{data set 3}
\caption{Data set 3: positive (blue dots) and negative (red crosses) training points.}
\label{fig:ds3}
\end{figure}

An RBF SVM with $\gamma=100$ and $C=1$ gives the curved boundary of Fig.~\ref{fig:rbf100}, with a training error of 0.0569 (12 points)
and 126 support vectors. Fig.~\ref{fig:gamma} and Table~\ref{tab:gamma} compare $\gamma=1,10,100$ and $1000$ with $C=1$ fixed.
For $\gamma=1$ each kernel is wider than the whole data set, so the boundary is almost straight and the model underfits (error 0.0806). For
$\gamma=10$ the boundary starts to follow the data. For $\gamma=100$ it becomes wavy and traces the local shape of the class border.
For $\gamma=1000$ every kernel covers only its immediate neighbourhood, the boundary breaks into small islands, 208 of the 211 points are
support vectors and the training error is zero, which is clear overfitting.

\begin{figure}[!t]
\centering
\projfig{fig_rbf100.png}{0.8\columnwidth}{RBF boundary}
\caption{Decision boundary of the RBF SVM with $\gamma=100$ and $C=1$ (training error 0.0569).}
\label{fig:rbf100}
\end{figure}

\begin{figure*}[!t]
\centering
\projfig{fig_gamma.png}{0.9\textwidth}{boundaries for different gamma}
\caption{Decision boundaries of the RBF SVM ($C=1$) for $\gamma=1,10,100$ and $1000$ (left to right).}
\label{fig:gamma}
\end{figure*}

\begin{table}[!t]
\caption{RBF Kernel SVM on Data Set 3 ($C=1$)}
\label{tab:gamma}
\centering
\begin{tabular}{@{}rccc@{}}
\toprule
$\gamma$ & Training error & Misclassified & SV\\
\midrule
1 & 0.0806 & 17 & 94\\
10 & 0.0664 & 14 & 73\\
100 & 0.0569 & 12 & 126\\
1000 & 0.0000 & 0 & 208\\
\bottomrule
\end{tabular}
\end{table}

% =====================================================================
\section{Discussion}
The experiments show the two hyperparameters that control SVM complexity. In the linear case, $C$ trades margin width against training
errors. A small $C$ gives a wide margin defined by many support vectors, but a value that is too small underfits (data set 2 at $C=10^{-6}$).
A large $C$ reduces to the hard-margin SVM, whose boundary rests on a few points. Data sets 1 and 2 show that the useful range of $C$ depends on how
close the classes are, and that increasing $C$ beyond the largest dual variable of the hard-margin solution has no effect.

On the digits every tested model performs very well: the test error never exceeds 0.19\%, and for $C\ge3\times10^{-4}$ it is at most 0.095\%.
The hard-margin SVM fits the training set perfectly but is not the best on the test set. A little slack ($C\approx0.01$--$0.03$) halves the test
error from two images to one. These differences amount to one or two images out of 2115, so they are not statistically strong, but cross-validation
nevertheless picks a value from the good region without using the test set.

The misclassified images expose the limitation of a linear model on raw pixels: it matches ink against a fixed template, so very thick or
very slanted 1s, hooked or broken 0s and images with stray marks end up near or across the boundary.

For the RBF kernel, $\gamma$ plays a role similar to $C$. A small $\gamma$ gives a smooth, nearly linear boundary that underfits, and
a large $\gamma$ memorises the training set. Since the training error always favours the largest $\gamma$, in practice $\gamma$ and $C$ should
be selected by cross-validation; for data set 3 values around 10--100 look like a reasonable compromise.

% =====================================================================
\section{Conclusion}
We trained linear and RBF-kernel SVMs on two-dimensional data and on handwritten digits. On the two separable 2-D data sets
the SVM reached zero test error; lowering $C$ widened the margin and increased the number of support vectors, and a very small $C$ caused
underfitting. On the 0/1 digits the hard-margin SVM had zero training error and a test error of 0.095\%, and a moderately regularised model
($C=0.03$, chosen by five-fold cross-validation) reduced the test error to 0.047\%. Remaining errors come from atypical writing styles that a
linear pixel template cannot capture. On a non-separable 2-D data set the RBF kernel produced non-linear boundaries whose complexity grows with
$\gamma$, from underfitting at $\gamma=1$ to overfitting at $\gamma=1000$.

% =====================================================================
\FloatBarrier
\section*{Code Availability}
The code was written and run in Google Colab and is in the submitted Jupyter notebook \texttt{SVM\_Course\_Project.ipynb}. It uses NumPy, Matplotlib and scikit-learn.
The data files are uploaded to the Colab session at the start of the notebook.

% =====================================================================
\bibliographystyle{IEEEtran}
\bibliography{references}

\end{document}
